\subsubsection{A Priori Bounds}\label{subsec:elliptic-a-priori-bounds}

The maximum principle can be used to obtain pointwise estimates for solutions of $Lu=f$ in bounded domains. To state these bounds, we introduce the following quantities:
\begin{equation}\label{eq:elliptic-pointwise-ellipticity-and-drift}
\lambda(\bm{x})=\min_{\bm{\xi}\in\mathbb{R}^n\setminus\{0\}}\dfrac{a_{ij}(\bm{x})\xi_i\xi_j}{|\bm{\xi}|^2},\quad\beta=\max_{\bm{x}\in\overline\Omega}\dfrac{|\bm{b}(\bm{x})|}{\lambda(\bm{x})}.
\end{equation}


We have the following result.

\begin{Theorem}
\label{thm:elliptic-strip-a-priori-bound}
Assume that $\Omega$ is bounded and contained in the strip between two parallel planes of distance $d$. Assume also that $c\leq0$. If $Lu=f$, then
\begin{equation}\label{eq:elliptic-strip-a-priori-bound}
\max_{\overline\Omega}|u|\leq\max_{\partial\Omega}|u|+C\max_{\overline\Omega}\dfrac{|f|}{\lambda}.
\end{equation}
Here $C=\exp((\beta+1)d)-1$.
\end{Theorem}

\begin{Proof}
We shall assume $Lu\geq f$ and prove that
\begin{equation}\label{eq:elliptic-strip-upper-bound}
\max_{\overline\Omega}u\leq\max_{\partial\Omega}u^++C\max_{\overline\Omega}\dfrac{|f^-|}{\lambda}.
\end{equation}
The theorem follows by applying this inequality to both $u$ and $-u$. Since $\lambda$ and $\beta$ do not change under rotations of the coordinate system, we may assume that $\Omega$ is contained in the strip $0<x_1<d$. We set $L_0=L-c$. For $\alpha\geq\beta+1$, we have
\begin{equation}\label{eq:elliptic-strip-exponential-barrier}
L_0\exp(\alpha x_1)=(\alpha^2a_{11}+\alpha b_1)\exp(\alpha x_1)\geq\lambda(\alpha^2-\alpha\beta)\exp(\alpha x_1)\geq\lambda.
\end{equation}
Let
\begin{equation}\label{eq:elliptic-strip-comparison-function}
v=\max_{\partial\Omega}u^++(\exp(\alpha d)-\exp(\alpha x_1))\max_{\overline\Omega}\dfrac{|f^-|}{\lambda}.
\end{equation}
Then obviously $v\geq u$ on $\partial\Omega$, and $v\geq0$ on $\Omega$. We compute
\begin{equation}\label{eq:elliptic-strip-comparison-inequality}
L(v-u)=L_0v+cv-Lu\leq-\lambda\max_{\overline\Omega}\dfrac{|f^-|}{\lambda}-f\leq0.
\end{equation}
Corollary~\ref{cor:elliptic-comparison} shows that $u\leq v$, which yields the desired result for $C=\exp(\alpha d)-1$.
\end{Proof}

The following corollary shows that, in certain cases, it is possible to dispense with the requirement that $c\leq0$.

\begin{Corollary}
\label{cor:elliptic-small-positive-zero-order}
Let $Lu=f$ in a bounded domain $\Omega$. Let $C$ be as in Theorem~\ref{thm:elliptic-strip-a-priori-bound} and assume that
\begin{equation}\label{eq:elliptic-positive-zero-order-smallness}
C^*=1-C\max_{\overline\Omega}\dfrac{c^+}{\lambda}>0.
\end{equation}
Then
\begin{equation}\label{eq:elliptic-small-positive-zero-order-bound}
\max_{\overline\Omega}|u|\leq\dfrac1{C^*}\left(\max_{\partial\Omega}|u|+C\max_{\overline\Omega}\dfrac{|f|}{\lambda}\right).
\end{equation}

\end{Corollary}

The proof follows by applying the result of the previous theorem to the equation $(L_0+c^-)u=f-c^+u$.

\paragraph*{Problems}

\begin{enumerate}[label=\thesection.\arabic*.,ref=\thesection.\arabic*,start=1,series=book-problems]

\item\label{ex:semilinear-cubic-maximum-bound} Let $u$ be a solution of $\Delta u=u^3-u$ on a bounded domain $\Omega$. Assume that $u=0$ on $\partial\Omega$. Show that $u\in[-1,1]$ throughout $\Omega$. Can the values $\pm1$ be achieved?

\item\label{ex:trace-positive-semidefinite-product} Assume that the $n\times n$ matrices $\bm{A}$ and $\bm{B}$ are symmetric and positive semidefinite. Show that $\operatorname{tr}(\bm{A}\bm{B})\geq0$. Hint: $\bm{B}=\sum_k\lambda_k\bm{q}_k\bm{q}_k^T$, where the $\lambda_k$ and $\bm{q}_k$ are the eigenvalues and eigenvectors of $\bm{B}$.

\item\label{ex:weak-maximum-positive-zero-order-counterexample} Give a counterexample showing that Corollary~\ref{cor:weak-maximum-nonpositive-zero-order} does not hold if $c>0$.

\item\label{ex:unbounded-elliptic-comparison-counterexample} Show that Corollary~\ref{cor:elliptic-comparison} fails if $\Omega$ is unbounded. Hint: Consider the problem $\Delta u=0$, $u=0$ on $\partial\Omega$ when $\Omega$ is a strip bounded by parallel planes.

\item\label{ex:interior-tangent-ball} If $\partial\Omega$ is of class $C^2$ and $\bm{x}_0$ is on $\partial\Omega$, show that there is a ball lying entirely in $\Omega$ with $\bm{x}_0$ on its boundary.

\item\label{ex:drift-neumann-adjoint-solvability} (a) On the bounded domain $\Omega$ with smooth boundary, let $u$ be a solution of the problem
\begin{equation}\label{eq:drift-neumann-problem}
\Delta u+a_i(\bm{x})\dfrac{\partial u}{\partial x_i}=f(\bm{x}),\quad\dfrac{\partial u}{\partial n}=0\text{ on }\partial\Omega.
\end{equation}
Assume that $f\geq0$ in $\Omega$. Show that $u$ is a constant and $f=0$.

(b) Show that problem \eqref{eq:drift-neumann-problem} can have a solution only if
\begin{equation}\label{eq:drift-neumann-adjoint-compatibility}
\int_\Omega f(\bm{x})v(\bm{x})\,\mathrm{d}\bm{x}=0
\end{equation}
for every solution $v$ of the "adjoint" equation
\begin{equation}\label{eq:drift-neumann-adjoint-equation}
\Delta v-\dfrac{\partial}{\partial x_i}\left(a_i(\bm{x})v\right)=0,\quad\dfrac{\partial v}{\partial n}-a_i(\bm{x})n_iv=0\text{ on }\partial\Omega.
\end{equation}
(c) Using techniques to be developed in later chapters, one can show that the condition \eqref{eq:drift-neumann-adjoint-compatibility} is also sufficient and that the solution space of \eqref{eq:drift-neumann-adjoint-equation} is one-dimensional. Taking these facts for granted, show that solutions of \eqref{eq:drift-neumann-adjoint-equation} are either non-negative or non-positive. Equations of the form \eqref{eq:drift-neumann-adjoint-equation} are called Fokker-Planck equations and arise in statistical physics. Only non-negative solutions are physically meaningful.

\item\label{ex:hexagon-laplace-eigenvalue-lower-bound} Let $\Omega$ be a regular hexagon with side $a$. Let $\lambda\in\mathbb{R}$ be such that the equation $\Delta u+\lambda u=0$ with boundary condition $u=0$ has a nontrivial solution in $\Omega$. Give a lower bound for $\lambda$.

\end{enumerate}
